% Ad Familiares 12.17 --- headnote
% ad-familiares-12-17, mid-September 46 BC, Rome

\textit{Cicero to Q. Cornificius, from Rome around the middle of September 46 BC ---
Perseus dateline \textit{Scr. Romae circ. med. in Sept. a. 708 (46)}.
Cornificius, addressed here as \textit{conlega} (a fellow-augur), has gone out
to a provincial command --- the troubles in Syria mentioned in section 1 (the
Caecilius Bassus affair) sit near his theatre --- and Cicero, finding Rome
quiet to the point of stagnation under Caesar's dictatorship, fills the
silence with literature. The book he wants Cornificius to read is the
\textit{Orator} (``de optimo genere dicendi''), addressed to Brutus and
finished in the same months; the joke about voting for it ``as a favour''
if the substance fails to convince is the relaxed shorthand of one
scholar-statesman to another. The closing tribute --- ``no one above you,
few beside you'' --- is the warmth Cicero reserves for the small circle of
educated younger men whose career he means to back.}
